\chapter{Introdução}\label{cap:introducao}
O estudo da Matemática no Ensino Básico brasileiro, essencial para o desenvolvimento do pensamento lógico e da capacidade de análise crítica, enfrenta muitos desafios, dentre eles, um dos mais persistentes: o desinteresse e a descontextualização da disciplina. Para muitos estudantes, a Matemática é percebida como um conjunto de fórmulas abstratas, o que contribui para o baixo desempenho e para a visão de que a matéria é ``chata'' e pouco relevante para a vida prática. Essa crítica à abordagem tradicional é defendida por \citeonline{dambrosio1999etnomatematica}, que argumenta que a ênfase excessiva em um modelo eurocêntrico e a-histórico resulta em uma ``matemática obsoleta e desinteressante'' para o aluno, tornando imprescindível que o ensino da matemática busque metodologias que resgatem o significado e a atratividade do saber.

A Geometria, um campo que exige a visualização espacial e a compreensão rigorosa de relações complexas, é particularmente vulnerável a essa crise. O abandono do ensino da Geometria no Brasil é um fato histórico e curricular amplamente documentado por \citeonline{pavanello1993abandono}. Segundo a autora, a falha reside na priorização de um ensino excessivamente formal e abstrato, onde a didática se concentra na memorização de propriedades e na apresentação dos teoremas exclusivamente de forma escrita e robusta. Esse modelo, ao não fornecer o suporte visual necessário, impede a transição do pensamento intuitivo para o pensamento dedutivo eficaz, transformando o rigor das demonstrações geométricas em um obstáculo intransponível. \citeonline{pavanello1993abandono}, também conclui que o trabalho com a Geometria, quando bem conduzido, deve favorecer o desenvolvimento de um pensamento crítico e autônomo. É nesse contexto de desengajamento e falha na transposição didática que a inovação metodológica se torna uma solução viável para resgatar o valor da disciplina.

%% ------- JUSTIFICATIVCA PARA O TEMA -- INICIO  

\section{Justificativa para a escolha do tema}

A escolha do tema desta dissertação, focado na Geometria Plana Olímpica e na produção multimídia, é justificada pela necessidade de intervir diretamente na crise de interesse, nas dificuldades de visualização e rigor matemático no Ensino Médio observado por \citeonline{pavanello1993abandono}. O presente projeto se fundamenta em três pilares essenciais: a filosofia da resolução de problemas, a progressão didática rigorosa e o imperativo da inovação tecnológica.

Primeiramente, a opção pelo formato olímpico visa elevar o nível de desafio e engajamento, priorizando o desenvolvimento do raciocínio sobre a memorização. A Geometria Olímpica, neste contexto, é utilizada como o veículo para o desenvolvimento da Heurística, essencial para que o estudante se torne um resolvedor de problemas autônomo. Essa postura é endossada por \citeonline{polya2006arte}, que enfatiza a importância da prática e da descoberta pessoal:

\begin{citacao}
Uma grande descoberta resolve um grande problema, mas há sempre uma pitada de descoberta na resolução de qualquer problema. Se ele desafiar a curiosidade e puser em jogo as faculdades inventivas, quem o resolver por seus próprios meios, experimentará a tensão e gozará o triunfo da descoberta \cite[p.~V]{polya2006arte}.
\end{citacao}

A trilha de estudos, portanto, será estruturada com base nas etapas do método de resolução de problemas, guiando o aluno na compreensão do problema, na elaboração de um plano de solução e na sua execução.

Em segundo lugar, a trilha aborda o problema da progressão didática. O rigor da Geometria exige que o aluno avance sequencialmente no seu entendimento das formas. Este trabalho adota a teoria de \citeonline{vanhiele1986structure} como base para a estruturação dos módulos, que postula que o desenvolvimento do pensamento geométrico ocorre em níveis hierárquicos. O material será desenhado para guiar o estudante dos Níveis Iniciais (Visualização) até o Nível de Dedução Formal, onde as demonstrações são compreendidas em sua essência lógica.

Por fim, o foco na produção de vídeos de animação e em uma plataforma digital é a justificativa metodológica para superar a falha do material escrito. Esta abordagem é a intervenção direta na crítica de Pavanello, fornecendo o suporte visual e dinâmico que o texto estático não consegue oferecer.


\section{Revisão Bibliográfica}
%%%% --------------------
O primeiro pilar desta revisão é o desenvolvimento do pensamento geométrico. O modelo de \citeonline{vanhiele1986structure} é a referência principal para o desenho didático da trilha de estudos, pois diagnostica a principal falha na transposição da Geometria no Ensino Básico: a incompatibilidade entre o nível de raciocínio do aluno e a complexidade do material apresentado. A teoria postula que o pensamento evolui através de cinco níveis hierárquicos, sendo impossível ao estudante ascender a um nível superior sem a completa assimilação do anterior. O rigor das demonstrações geométricas, que é o ponto focal da Geometria Olímpica, corresponde ao Nível 3 (Dedução Formal). O fracasso em atingir este nível é validado pela literatura nacional, que aponta a desconexão entre o ensino e o aprendizado. Conforme observaram \citeonline[p.3]{kaleff1994desenvolvimento}, ao investigarem o modelo no contexto brasileiro: ``Seus trabalhos revelam uma alarmante falta de harmonia entre o ensino e aprendizado em matemática''.

Para este trabalho, a teoria de \citeonline{vanhiele1986structure} é utilizada para guiar a sequencialidade da trilha. A trilha de estudos é, portanto, uma intervenção didática intencional que utilizará o recurso multimídia como o principal catalisador para a transição entre os níveis, do reconhecimento intuitivo até a prova rigorosa.

O segundo pilar é a Filosofia da Resolução de Problemas. A trilha utiliza a Geometria Olímpica como veículo para o desenvolvimento da Heurística, essencial para que o estudante se torne um resolvedor de problemas autônomo. O rigor é alcançado através do método de \citeonline{polya2006arte}, cujas etapas guiam o aluno na compreensão do problema, na elaboração de um plano de solução e na sua execução. O material didático da trilha será desenhado para incentivar ativamente esse processo de descoberta e invenção.

O terceiro pilar é o Desenho da Instrução Multimídia. A proposta de utilizar vídeos de 3 minutos, programados em JavaScript e integrados ao texto, é a resposta metodológica à falha do material estático. Esta metodologia é validada pela Teoria Cognitiva da Aprendizagem Multimídia (TCAM) de \citeonline{mayer2021multimedia}. O principal foco da TCAM é o gerenciamento da Capacidade Limitada da memória de trabalho do aluno, que tende a ser sobrecarregada por demonstrações complexas e puramente textuais. Este material visa aplicar diretamente três princípios centrais:

\begin{itemize}
    \item O Princípio Multimídia justifica a combinação do texto com a animação (imagem), pois as pessoas aprendem melhor quando se combina palavras e imagens do que apenas palavras \cite{mayer2021multimedia}.
    \item O Princípio da Contiguidade Temporal sustenta a exibição do texto ao lado do vídeo, garantindo que o aluno processe o que vê e o que ouve simultaneamente.

    \item O Princípio da Segmentação valida o formato de vídeos curtos (3 minutos), pois a divisão do conteúdo em passos menores e controláveis evita a sobrecarga cognitiva (Capacidade Limitada).
\end{itemize}


%%%%---------------------


% ------ OBJETIVOS 

\section{Objetivos}
\subsection{Objetivos Gerais}
O presente trabalho tem como objetivo geral desenvolver e documentar uma Trilha de Estudos em Geometria Plana Olímpica, com foco na produção de um Produto Educacional multimídia (combinando texto rigoroso e vídeos de animação), analisando seu potencial teórico-didático para promover o avanço do raciocínio dedutivo no Ensino Médio.
\subsection{Objetivos específicos}

Para a consecução do objetivo geral, esta dissertação propõe-se a:

\begin{enumerate}
    \item Reunir e formalizar as demonstrações e os exercícios dos principais teoremas e corolários de Geometria Plana, garantindo o rigor matemático necessário para o nível olímpico (Heurística de Pólya).
    
    \item Estruturar a progressão didática da Trilha em módulos sequenciais, utilizando a Teoria de \citeonline{vanhiele1986structure} para guiar o estudante do Nível de Visualização ao Nível de Dedução Formal.
    
    \item Desenvolver a metodologia de produção dos vídeos de animação (curtos, em JavaScript), detalhando o processo de storyboard e narração para a ilustração visual dos conceitos.
    
    \item Analisar o design metodológico do Produto Educacional à luz dos Princípios da Teoria Cognitiva da Aprendizagem Multimídia (TCAM) de \citeonline{mayer2021multimedia}, justificando a eficácia da combinação texto-vídeo para gerenciar a carga cognitiva do aluno.
    
    \item Apresentar o Produto Educacional completo (Trilha de Estudos e Vídeos) de forma documentada e organizada, pronto para futura aplicação em ambiente de Ensino Médio.
\end{enumerate}

 %%% ------- PROCEDIMENTOS METODOLOGICOS
 
\section{Procedimentos metodológicos}
A presente dissertação configura-se como uma Pesquisa Bibliográfica e de Desenvolvimento Teórico, cujo foco principal é a criação e a fundamentação de um Produto Educacional (a Trilha Multimídia) para intervir em um problema didático da prática.

\subsection{Abordagem e Natureza}
\begin{itemize}
    \item Natureza da Pesquisa: O trabalho possui a natureza de Pesquisa Bibliográfica de Desenvolvimento Teórico. Ele se alinha à definição de Pesquisa Exploratória de \citeonline{gerhardt2009metodos}, que é realizada quando o objetivo é proporcionar maior familiaridade com um problema, buscando referências para a criação de um modelo teórico e didático.
    \item Abordagem: O trabalho adota uma abordagem qualitativa, uma vez que visa à análise interpretativa das características, do design e das potencialidades do Produto Educacional desenvolvido. Essa escolha metodológica se justifica porque, conforme \citeonline{ludke1986pesquisa}, na pesquisa qualitativa o pesquisador é o principal instrumento de investigação, atuando diretamente na interpretação do objeto em seu contexto, o que é fundamental para alcançar os objetivos propostos.
\end{itemize}

\subsection{Objetivos e Procedimentos}
\begin{itemize}
    \item Objetivos da Pesquisa: O estudo é de natureza Exploratório-Descritiva. É exploratório ao buscar na literatura a crítica ao ensino da Geometria e as soluções multimídia. É descritivo ao detalhar a estrutura, o conteúdo e a metodologia de design da Trilha de Estudos e dos vídeos de animação \cite{gerhardt2009metodos}.
    \item Procedimentos Técnicos: O trabalho baseia-se primordialmente na Pesquisa Bibliográfica. Foram analisadas obras e artigos de referência nos campos da Didática da Matemática, da Tecnologia Educacional e da historiografia do ensino da Geometria. Esta pesquisa embasa todas as decisões tomadas no design do produto.
\end{itemize}

%%% ----- PRODUTO EDUCACIONAL

\section{Produto Educacional}

A culminação desta pesquisa bibliográfica e metodológica é o Produto Educacional, que se materializa na forma de uma Trilha de Estudos Digital em Geometria Plana Olímpica. Este recurso, desenvolvido sob as diretrizes do Mestrado Profissional, visa ser um artefato didático de alta replicabilidade e eficácia teórica. \\
O Recurso Educacional é composto por:

\begin{enumerate}
    \item Material Escrito (Texto e Estrutura): Consiste nas demonstrações rigorosas dos teoremas e nas listas de exercícios resolvidos, organizados sequencialmente conforme os Níveis de Pensamento Geométrico de \citeonline{vanhiele1986structure}.

\item Material Audiovisual: É a parte central do projeto. Consiste em uma série de vídeos curtos de animação (cerca de 3 minutos), estilo Reels/TikTok, que ilustram as etapas mais complexas das demonstrações geométricas. Conforme a Figura \ref{fig:flame-de-animacao}, que ilustra um frame da animação, o recurso é desenhado para guiar visualmente o aluno através das construções auxiliares necessárias para a prova dos teoremas.
\begin{figure}[!h]
    \caption{flames de animação: (a) bissetriz interna, (b) ENQ 2025} 
    \vspace{0.25cm}
	\label{fig:flame-de-animacao}
	\centering
	\subfloat[]{		\includegraphics[scale=0.13]{figuras/print1.jpg}
	}\hspace{0.2cm} 
	\subfloat[]{
		\includegraphics[scale=0.13]{figuras/print2.jpg}
	}	
	\fonte{Autoria própria (2025)}
\end{figure}

\end{enumerate}

A produção deste Material Audiovisual é realizada através de um software de animação especializado, que agiliza a criação do conteúdo multimídia. Esta escolha metodológica garante a alta qualidade visual, a rapidez na produção do material e a capacidade de focar o esforço do pesquisador no design didático, e não na codificação bruta.

A Trilha de Estudos completa será disponibilizada em uma plataforma digital para facilitar o acesso e promover a integração entre o texto e os vídeos de suporte, conforme exigido pelos Princípios de Contiguidade de \citeonline{mayer2021multimedia}.



\section{Estrutura da Dissertação}
Este trabalho de dissertação está organizado em cinco capítulos, além do Produto Educacional que é anexo e peça central da pesquisa. A divisão está estabelecida da seguinte forma:

\begin{itemize}
    \item Capítulo 1 – Introdução: Apresenta o contexto temático da crise de interesse na Matemática, o problema de pesquisa, os objetivos (geral e específicos), a relevância e a justificativa para a escolha do tema (Geometria, Olimpíadas, Multimídia).
    \item Capítulo 2 – Fundamentação Teórica: Detalha os pilares conceituais que sustentam a intervenção didática. Serão exploradas a teoria da Resolução de Problemas (George Pólya), a teoria de Desenvolvimento do Pensamento Geométrico (Van Hiele), e a Teoria Cognitiva de Aprendizagem Multimídia (TCAM) (Richard Mayer), que justifica o uso dos vídeos de animação.
    \item Capítulo 3 – Metodologia: Descreve o design da pesquisa, que é de natureza teórica, detalhando o processo de elaboração do Produto Educacional, a definição do público-alvo e os procedimentos para a criação das animações em JavaScript.
    \item Capítulo 4 – O Produto Educacional: Apresenta, em detalhes, a Trilha Multimídia de Geometria Plana Olímpica. Este capítulo contém o cerne do trabalho, incluindo as demonstrações formais, a progressão dos exercícios resolvidos e a análise teórica da eficácia didática do material.
    \item Capítulo 5 – Conclusões e Sugestões para Trabalhos Futuros: Sintetiza os principais achados da pesquisa, verifica o cumprimento dos objetivos propostos e sugere caminhos para futuras investigações e desdobramentos do Produto Educacional.
\end{itemize}

O Produto Educacional, composto pela Trilha Multimídia de Geometria Plana Olímpica e seus vídeos de animação, é o anexo deste documento e a materialização desta pesquisa.

 

   % ----- CRONOGRAMA DA DISSERTAÇÃO
   
\section{Cronograma da Dissertação PROFMAT: Trilha de Geometria Multimídia}
\begin{itemize}
    \item Prazo Total: 12 Meses (Jan/2026 a Dez/2026)
    \item Entrega à Banca (Leitura): 01 de Novembro de 2026
    \item Defesa Prevista: Semana de 07 a 11 de Dezembro de 2026
\end{itemize}

O cronograma da dissertação está dividido em quatro fases estratégicas, com o objetivo de otimizar o tempo e garantir o foco na produção do Produto Educacional. A Fase 1 (tabela \ref{tab:cronograma-fase1}) concentra-se na leitura e no embasamento teórico (Capítulos 2 e 3). A Fase 2 (tabela \ref{tab:cronograma-fase2}) é dedicada à criação e produção multimídia do Produto Educacional (Capítulo 4). A Fase 3 (tabela \ref{tab:cronograma-fase3}) abrange a escrita dos capítulos finais e a revisão formal do texto. Por fim, a Fase 4 (tabela \ref{tab:cronograma-fase4}) é o período destinado à entrega e defesa final.

%% CRONOGRAMA DISSERTAÇÃO PARTE 1

\begin{table}[!ht]
    \centering
    \caption{Cronograma da Dissertação - Fase 1: Fundação e Embasamento Teórico}
    
    \begin{tabular}{p{2cm} p{2.5cm} p{7cm} p{2.5cm}}
    \hline
    \textbf{Mês} & \textbf{Foco Principal} & \textbf{Tarefas Detalhadas} & \textbf{Entregável/Meta}\\
    \hline
    JAN/2026 & Planejamento Inicial & 1. Reunião de alinhamento com o Orientador. 2. Definição final da lista de Teoremas (a Trilha). 3. Busca e coleta de referências primárias (Pólya, Mayer, Van Hiele). & Fichamento das Referências.\\
    \hline
    FEV/2026 & Fundamentação Teórica I & 1. Leitura aprofundada e organização das citações. 2. Escrita da seção sobre Geometria, Polya e Van Hiele. & Capítulo 2 (Fundamentação Teórica) - 1ª Versão Bruta.\\
    \hline
    MAR/2026 & Fundamentação Teórica II & 1. Finalização do Capítulo 2 (Foco na TCAM de Mayer). 2. Definição e escrita da Metodologia de Design da Trilha. & Capítulo 3 (Metodologia de Design) - 1ª Versão Pronta.\\
    \hline
    \end{tabular} \label{tab:cronograma-fase1}
    \fonte{Autoria Própria (2025)}

\end{table}


% FAASE DOIS DO CRONOGRAMA

\begin{table}[!ht]
    \centering
    \caption{Cronograma da Dissertação - Fase 2: Design e Produção do Produto}
    
    \begin{tabular}{p{2cm} p{2.7cm} p{7cm} p{2.5cm}}
    \hline
    \textbf{Mês} & \textbf{Foco Principal} & \textbf{Tarefas Detalhadas} & \textbf{Entregável/Meta}\\
    \hline
    ABR/2026 & Roteirização do Produto & 1. Elaboração do Roteiro Didático da Trilha (Sequência Lógica de Teoremas). 2. Roteirização e storyboard detalhado dos primeiros vídeos de animação. & Mapa da Trilha (Sumário Detalhado do Cap. 4).\\
    \hline
    MAI/2026 & Produção do Conteúdo (Texto e LaTeX) & 1. Criação das demonstrações rigorosas em LaTeX para o material escrito. 2. Início da produção das animações em JavaScript (Primeiros 50\% dos vídeos). & 50\% das Demonstrações Escritas (Cap. 4) finalizadas.\\
    \hline
    JUN/2026 & Finalização da Produção Multimídia & 1. Finalização das demonstrações em LaTeX e dos exercícios resolvidos. 2. Finalização da produção e narração dos vídeos restantes. & Capítulo 4 (Produto Educacional) - Versão de Conteúdo (100\% do material criado).\\
    \hline
    JUL/2026 & Integração e Revisão do Produto & 1. Montagem final do Produto Educacional (Integração dos vídeos e textos no site/plataforma). 2. Revisão inicial do Orientador sobre o Produto. & Produto Educacional (Trilha e Vídeos) 100\% Concluído.\\
    \hline
    \end{tabular} \label{tab:cronograma-fase2}
    \fonte{Autoria Própria (2025)}
\end{table}

% FASSE 3 DO CRONOGRAMA

\begin{table}[!ht]
    \centering
    \caption{Cronograma da Dissertação - Fase 3: Escrita Acadêmica Final e Revisão}
    
    \begin{tabular}{p{2cm} p{2.5cm} p{7cm} p{2.5cm}}
    \hline
    \textbf{Mês} & \textbf{Foco Principal} & \textbf{Tarefas Detalhadas} & \textbf{Entregável/Meta}\\
    \hline
    AGO/2026 & Escrita Acadêmica Final & 1. Escrita do Capítulo 4 (Detalhando o Produto e sua Análise Teórica). 2. Escrita do Capítulo 5 (Considerações Finais e Contribuições). 3. Escrita Final do Capítulo 1 (Introdução). & Versão Bruta dos Capítulos 1, 4 e 5.\\
    \hline
    SET/2026 & Revisão e Formatação & 1. Revisão completa de linguagem, coesão e argumentação. 2. Formatação ABNT (Referências, Citações e Elementos Pré/Pós-Textuais). & Dissertação (Versão de Revisão) Pronta.\\
    \hline
    OUT/2026 & Aprovação Final do Orientador & 1. Revisão final minuciosa do Orientador. 2. Incorporação dos últimos ajustes. & 30 OUT: Dissertação (Versão Final de Entrega) APROVADA pelo Orientador.\\
    \hline
    \end{tabular} \label{tab:cronograma-fase3}
    \fonte{Autoria Própria (2025)}
\end{table}

%FASE FINAL DO CRONOGRAMA

\begin{table}[!ht]
    \centering
    \caption{Cronograma da Dissertação - Fase 4: Entrega e Defesa}
    
    \begin{tabular}{p{2cm} p{3cm} p{6.5cm} p{2.5cm}}
    \hline
    \textbf{Mês} & \textbf{Foco Principal} & \textbf{Tarefas Detalhadas} & \textbf{Entregável/Meta}\\
    \hline
    NOV/2026 & Entrega à Banca & 1. Finalização da impressão e encadernação (se necessário). 2. Cumprimento dos requisitos burocráticos (documentos, cópias). & 01 NOV: Entrega obrigatória da versão final à banca para leitura \\
    \hline
    DEZ/2026 & Defesa & 1. Preparação da apresentação de defesa (slides). 2. Simulação da defesa. & Semana de 07 a 11 DEZ:Defesa Publica da dissertação. \\
    \hline
    \end{tabular} \label{tab:cronograma-fase4}
    \fonte{Autoria Própria (2025)}
\end{table}
